\documentclass[11pt]{article}
\usepackage[margin=1in]{geometry}
\usepackage{amsmath, amssymb}
\usepackage{enumitem}
\usepackage{titlesec}
\usepackage{hyperref}
\usepackage{booktabs}
\usepackage{xcolor}
\usepackage{parskip}

\hypersetup{colorlinks=true, linkcolor=blue, urlcolor=blue}

\title{\textbf{EE627 --- Class Summary}\\[0.4em]
       \large March 13, 2026}
\author{Prof.\ Rensheng Wang}
\date{}

\begin{document}
\maketitle
\tableofcontents
\newpage

%------------------------------------------------------------
\section{Why Feature Selection?}
%------------------------------------------------------------

A common question: \textit{why not give all features to the ML algorithm
and let it decide which ones matter?} Three reasons argue against this:

\subsection{Reason 1: Overfitting}
When the number of features (columns) exceeds the number of samples
(rows), the model can fit the training data perfectly --- but it learns
both the true signal \emph{and} the noise.
\begin{itemize}
    \item Result: near-perfect training accuracy, terrible test accuracy.
    \item The model has memorized the noise features as if they were signal.
\end{itemize}

\subsection{Reason 2: Occam's Razor}
Simpler models are more explainable. A model with hundreds of features
is a black box even if it is technically more accurate. Interpretability
matters both for debugging and for stakeholder communication.

\subsection{Reason 3: Garbage In, Garbage Out}
Most real datasets contain many non-informative features.
\begin{itemize}
    \item Poor-quality inputs distract the learning process and degrade
      output quality.
    \item More features $\Rightarrow$ higher training cost, longer runtime,
      and harder implementation.
\end{itemize}

\textbf{Conclusion:} select only the useful features.


\newpage
%------------------------------------------------------------
\section{Three Families of Feature Selection Methods}
%------------------------------------------------------------

\begin{center}
\begin{tabular}{@{}lll@{}}
\toprule
\textbf{Family} & \textbf{How it works} & \textbf{Examples} \\
\midrule
Filter-based  & Score each feature independently; keep top $n$ & Pearson correlation, $\chi^2$ \\
Wrapper-based & Treat selection as a search problem over subsets & RFE, sequential forward selection \\
Embedded      & Algorithm has built-in selection during training & Lasso, Random Forest \\
\bottomrule
\end{tabular}
\end{center}


\newpage
%------------------------------------------------------------
\section{Filter Method: Pearson Correlation}
%------------------------------------------------------------

\subsection{Definition}
Pearson correlation between feature $x$ and target $y$:
\[
r = \frac{\sum (x - \bar{x})(y - \bar{y})}
         {\sqrt{\sum (x - \bar{x})^2 \sum (y - \bar{y})^2}}
\]
where $\bar{x}$ is the sample mean of $x$.

\subsection{Application}
Compute $|r|$ between each numerical feature and the target.
Keep the top $n$ features ranked by $|r|$.

\subsection{Python Implementation}
\begin{verbatim}
def cor_selector(X, y, num_feats):
    cor_list = []
    feature_name = X.columns.tolist()
    for i in X.columns.tolist():
        cor = np.corrcoef(X[i], y)[0, 1]
        cor_list.append(cor)
    cor_list = [0 if np.isnan(i) else i for i in cor_list]
    cor_feature = X.iloc[:,
        np.argsort(np.abs(cor_list))[-num_feats:]].columns.tolist()
    cor_support = [True if i in cor_feature else False
                   for i in feature_name]
    return cor_support, cor_feature

cor_support, cor_feature = cor_selector(X, y, num_feats)
\end{verbatim}


\newpage
%------------------------------------------------------------
\section{Wrapper Method: Recursive Feature Elimination (RFE)}
%------------------------------------------------------------

\subsection{Algorithm}
RFE is a \textbf{wrapper method}: it frames feature selection as a
search problem.

\begin{enumerate}
    \item Train the estimator on the current full feature set.
    \item Obtain the importance of each feature via
      \texttt{coef\_attribute} or \texttt{feature\_importances\_attribute}.
    \item Prune the least important features.
    \item Repeat on the reduced set until the desired number of
      features remains.
\end{enumerate}

\subsection{Python Implementation (with Logistic Regression)}
\begin{verbatim}
from sklearn.feature_selection import RFE
from sklearn.linear_model import LogisticRegression

rfe_selector = RFE(estimator=LogisticRegression(),
                   n_features_to_select=num_feats,
                   step=10, verbose=5)
rfe_selector.fit(X_norm, y)
rfe_support = rfe_selector.get_support()
rfe_feature = X.loc[:, rfe_support].columns.tolist()
print(str(len(rfe_feature)), 'selected features')
\end{verbatim}

The \texttt{step=10} parameter removes 10 features per iteration.
The \texttt{LogisticRegression} estimator provides \texttt{coef\_}
as the importance measure.


\newpage
%------------------------------------------------------------
\section{Embedded Methods: Lasso and Ridge Regression}
%------------------------------------------------------------

\subsection{Overfitting and Regularization}
Linear regression minimizes the residual sum of squares (RSS):
\[
\sum_{i=1}^{M} (y_i - \hat{y}_i)^2 =
\sum_{i=1}^{M}\!\left(y_i - \sum_{j=0}^{n} w_j x_{ij} - b\right)^{\!2}
\]
With many features and limited data, this leads to overfitting.
\textbf{Regularization} adds a penalty term to shrink or zero out
weights.

\subsection{Lasso (L1 Regularization)}
\[
\text{Cost}_{\text{Lasso}} =
\sum_{i=1}^{M}\!\left(y_i - \sum_{j=0}^{n} w_j x_{ij} - b\right)^{\!2}
+ \lambda \sum_{j=0}^{n} |w_j|
\]
\begin{itemize}
    \item The L1 penalty forces many weights \emph{exactly} to zero.
    \item Higher $\lambda$ $\Rightarrow$ fewer non-zero features.
    \item Built-in feature selection: features with zero weight
      are effectively eliminated.
\end{itemize}

\subsection{Ridge (L2 Regularization)}
\[
\text{Cost}_{\text{Ridge}} =
\sum_{i=1}^{M}\!\left(y_i - \sum_{j=0}^{n} w_j x_{ij} - b\right)^{\!2}
+ \lambda \sum_{j=0}^{n} w_j^2
\]
\begin{itemize}
    \item The L2 penalty shrinks weights toward zero but rarely
      sets them \emph{exactly} to zero.
    \item Ridge reduces complexity but does not eliminate features.
    \item Lasso is preferred when explicit feature selection is needed.
\end{itemize}

\subsection{Lasso: \texttt{SelectFromModel} in Python}
\begin{verbatim}
from sklearn.feature_selection import SelectFromModel
from sklearn.linear_model import LogisticRegression

embeded_lr_selector = SelectFromModel(
    LogisticRegression(penalty="l1"),
    max_features=num_feats)
embeded_lr_selector.fit(X_norm, y)
embeded_lr_support = embeded_lr_selector.get_support()
embeded_lr_feature = X.loc[:, embeded_lr_support].columns.tolist()
\end{verbatim}


\newpage
%------------------------------------------------------------
\section{Embedded Method: Tree-Based Feature Importance}
%------------------------------------------------------------

\subsection{How Random Forest Selects Features}
Random Forest (an ensemble of decision trees) provides feature
importance scores automatically:
\begin{itemize}
    \item Feature importance = mean improvement in the splitting
      criterion (``impurity'') produced by that variable across
      all trees.
    \item In other words: how much does splitting on this feature
      reduce impurity on average?
    \item Features can be ranked by importance and a subset selected.
\end{itemize}

\subsection{Python Implementation}
\begin{verbatim}
from sklearn.feature_selection import SelectFromModel
from sklearn.ensemble import RandomForestClassifier

embeded_rf_selector = SelectFromModel(
    RandomForestClassifier(n_estimators=100),
    max_features=num_feats)
embeded_rf_selector.fit(X, y)
embeded_rf_support = embeded_rf_selector.get_support()
embeded_rf_feature = X.loc[:, embeded_rf_support].columns.tolist()
\end{verbatim}


\newpage
%------------------------------------------------------------
\section{MATLAB Implementation: Sequential Forward Selection}
%------------------------------------------------------------

The course project uses MATLAB for feature selection via
\texttt{sequentialfs}, which performs \textbf{sequential forward
selection} (a wrapper method).

\subsection{Criterion Function}
A criterion function is required to score each candidate feature
subset. Here the criterion is the \textbf{deviance} from a binomial
logistic regression fit:

\begin{verbatim}
function dev = critfun(X, Y)
    [b, dev] = glmfit(X, Y > 0, 'binomial');
\end{verbatim}

Lower deviance = better fit = better feature subset.

\subsection{Feature Selection Driver}
\begin{verbatim}
function [feaInd] = featureSelection(TrainMat, target)
    maxdev = chi2inv(.95, 1);
    opt = statset('display', 'iter', ...
                  'TolFun', maxdev, ...
                  'TolTypeFun', 'abs');
    feaInd = sequentialfs(@critfun, TrainMat, target, ...
                          'cv', 'none', ...
                          'nullmodel', true, ...
                          'options', opt, ...
                          'direction', 'forward');
\end{verbatim}

Key parameters:
\begin{itemize}
    \item \texttt{direction = 'forward'}: start from empty set, add
      one feature at a time.
    \item \texttt{TolFun = chi2inv(.95,1)}: stop adding features when
      the improvement in deviance falls below the 95th percentile of
      the $\chi^2(1)$ distribution (i.e., $\approx 3.84$).
    \item \texttt{nullmodel = true}: compare each subset against an
      intercept-only model.
    \item \texttt{cv = 'none'}: no cross-validation; use training data
      directly.
\end{itemize}


\newpage
%------------------------------------------------------------
\section{Best Subset Selection with \texttt{bestglm}}
%------------------------------------------------------------

\subsection{The Best Subset Problem}
Given $p$ features, there are $2^p$ possible subsets. Best subset
selection exhaustively evaluates all of them and picks the one that
minimizes a chosen criterion.

\begin{itemize}
    \item The \texttt{bestglm} R package implements this for the GLM
      family (linear, logistic, etc.).
    \item Feasible in practice for $p \lesssim 15$--25; for larger
      $p$, stepwise methods are used instead.
    \item Input format: a data frame \texttt{Xy} where the last column
      is the response and all other columns are predictors.
\end{itemize}

\subsection{Information Criteria for Model Selection}
All criteria penalize model fit (deviance $\mathcal{D} = -2\log\hat{\mathcal{L}}$)
by a term that grows with the number of parameters $k$:

\begin{center}
\begin{tabular}{@{}lll@{}}
\toprule
\textbf{Criterion} & \textbf{Formula} & \textbf{Tendency} \\
\midrule
AIC  & $\mathcal{D} + 2k$ & More variables; asymptotically efficient \\
BIC  & $\mathcal{D} + k\log n$ & Fewer variables; consistent \\
BIC$_\gamma$ & $\mathcal{D} + k\log n + 2\gamma\log\binom{p}{k}$ & Adjustable sparsity \\
BIC$_q$  & $\mathcal{D} + k\log n - 2k\log\tfrac{q}{1-q}$ & Bernoulli prior on inclusion \\
\bottomrule
\end{tabular}
\end{center}

\begin{itemize}
    \item When $n > 7$, BIC penalty exceeds AIC, so BIC always selects
      $\leq$ as many features as AIC.
    \item BIC is the default in \texttt{bestglm}.
    \item BIC$_q$ with $q = 1/2$ reduces to BIC.
\end{itemize}

\subsection{Prostate Cancer Example}
\begin{verbatim}
library(bestglm)
data(zprostate)
train <- (zprostate[zprostate[,10],])[,-10]
X <- train[, 1:8]
y <- train[, 9]
Xy <- cbind(as.data.frame(X), lpsa = y)

bestglm(Xy, IC = "AIC")   # selects 7 variables
bestglm(Xy, IC = "BIC")   # selects 2: lcavol, lweight
\end{verbatim}

AIC selects 7 variables (including two that are not significant at
5\%), while BIC selects only 2 (the two strongest predictors).
This illustrates BIC's stronger preference for parsimony.

\subsection{Cross-Validation Methods in \texttt{bestglm}}
When \texttt{IC = "CV"} is used, the package selects the best model
via cross-validation:

\begin{center}
\begin{tabular}{@{}lll@{}}
\toprule
\textbf{Method} & \textbf{Description} & \textbf{Notes} \\
\midrule
Delete-$d$ & Random subsets of size $d$ as validation set & Recommended default \\
$K$-fold   & Data partitioned into $K$ folds; each fold serves as validation once & Use 1-SD rule for stability \\
LOOCV      & Each observation left out once & High variance; asymptotically $\equiv$ AIC \\
\bottomrule
\end{tabular}
\end{center}

The \textbf{one-standard-deviation rule} for $K$-fold CV: instead of
the model with the single lowest CV error, choose the most parsimonious
model whose CV error is within one standard deviation of the minimum.
This produces more stable, simpler models.

\end{document}
